\section{Transporting a finite trace certificate to detection growth}
\label{sec:trace-detection}

This subsection converts the finite trace defect into a lower bound for
the global detection thresholds.  Write $h_{m,1},\ldots,h_{m,m}$ for the
zeros of the monic probabilists' Hermite polynomial $\operatorname{He}_m$,
and put $D_m=\operatorname{diag}(h_{m,1},\ldots,h_{m,m})$.  Its spectrum
is symmetric, $\operatorname{tr}D_m=0$, and
$\operatorname{tr}D_m^2=m(m-1)$.  For a real diagonal matrix $X$, let
$R(X)$ be the difference between its largest and smallest entries and set
\[
 \mathcal O_m(X)=\int_{U(m)}e^{i\operatorname{tr}(D_mUXU^*)}\,dU,
 \qquad
 \mathfrak r_m=\inf\{R(X):\mathcal O_m(X)<0\},
\]
where Haar measure has mass one.  The symmetry of $D_m$ makes
$\mathcal O_m(X)$ real.

\begin{proposition}[An orbital radius bound]
\label{prop:orbital-unimodal-bound}
For every $m\ge2$,
\begin{equation}
 \mathfrak r_m^2\ge\frac{4\pi^2(m+1)}{3m^2},
 \qquad (m-1)\mathfrak r_m^2\ge\pi^2.
 \label{eq:orbital-unimodal-radius}
\end{equation}
The constant $\pi^2$ in the second inequality is attained for $m=2$.
\end{proposition}

\begin{proof}
We first prove the elementary inequality
\begin{equation}
 \frac{\sin u}{u}\ge\frac12-\frac{u^2}{2\pi^2}
 \quad (u\in\mathbb R),
 \label{eq:sinc-parabola}
\end{equation}
with the continuous value at zero.  The difference $g(u)$ between the
two sides is even.  On $[0,\pi]$, write
$h(u)=\sin u-u\cos u-u^3/\pi^2$.  Then
$g'(u)=-h(u)/u^2$ and
$h'(u)=u^2(\sin u/u-3/\pi^2)$.  The function $\sin u/u$ decreases
strictly from one to zero on this interval, since the derivative of
$u\cos u-\sin u$ is $-u\sin u$.  Thus $h$ first increases and then
decreases.  Its values at both endpoints are zero, so $h\ge0$.
Consequently $g$ decreases to $g(\pi)=0$.  On $[\pi,3\pi/2]$,
\[
 g''(u)=\left(-\frac1u+\frac2{u^3}\right)\sin u
              -\frac{2\cos u}{u^2}+\frac1{\pi^2}>0,
\]
and $g'(\pi)=0$.  Finally, for $u\ge3\pi/2$,
$g'(u)\ge u/\pi^2-1/u-1/u^2>0$.  The last inequality follows because
$u^3/\pi^2-u-1$ increases there and its initial value is
$15\pi/8-1>0$.  This proves \eqref{eq:sinc-parabola}.

Let $Y$ have an even probability density $p$ which is nonincreasing on
the positive half-line.  The Stieltjes measure
$\nu(da)=-2a\,dp(a)$ represents this law as a mixture of uniform laws on
$[-a,a]$; explicitly,
\[
 p(y)=\int_{a\ge |y|}\frac{\nu(da)}{2a},\qquad
 \int\nu(da)=1,
 \qquad
 \mathbb E Y^2=\frac13\int a^2\nu(da).
\]
These identities follow by Stieltjes integration by parts, first for
step functions and then by monotone approximation.  An atom at zero
can be included as a uniform law with $a=0$.  Integrating
\eqref{eq:sinc-parabola} therefore gives
\begin{equation}
 \mathbb E\cos Y\ge\frac12-\frac{3\mathbb E Y^2}{2\pi^2}.
 \label{eq:unimodal-cosine-bound}
\end{equation}

We apply this to $Y=\operatorname{tr}(D_mUXU^*)$.  We give the exact
map that makes its density unimodal.  For distinct entries
$x_1<\cdots<x_m$ of $X$, let $y^{(k)}_1\le\cdots\le y^{(k)}_k$ be the
eigenvalues of the upper left $k\times k$ corner of $UXU^*$.
Olshanski's corner-measure formula and its iterated form
\cite[Proposition 3.1 and Corollary 3.2, pp.~6--7]{olshanski2013}
send Haar measure to uniform Lebesgue measure on the convex
Gelfand--Tsetlin polytope
\[
 \{(y^{(1)},\ldots,y^{(m-1)}):
             y^{(1)}\prec\cdots\prec y^{(m-1)}\prec x\}.
\]
Here $\prec$ means interlacing.  If $s_k=\sum_jy^{(k)}_j$, $s_0=0$,
and $s_m=\operatorname{tr}X$, the $k$th diagonal entry of $UXU^*$ is
$s_k-s_{k-1}$.  Thus the map to $Y$ is precisely the linear functional
\[
 (y^{(1)},\ldots,y^{(m-1)})
       \longmapsto\sum_{k=1}^m h_{m,k}(s_k-s_{k-1}).
\]
For $m=2$ its density is uniform on an interval.  For $m>2$, the
Brunn--Minkowski inequality applied to parallel sections of this
convex polytope says that a positive power of the density is concave
on its support.  Conjugating $D_m$ by the permutation that reverses its
eigenvalues sends $Y$ to $-Y$ without changing its law.  The density is
therefore even, and concavity makes it nonincreasing on the positive
half-line.  Repeated entries of $X$ follow by approximation from
distinct entries: the orbital integrals and their second moments are
continuous in $X$.  Scalar $X$ gives $Y=0$ directly.

For completeness, its variance can be computed without an orbital
integration formula.  The squared absolute entries of one Haar row
are uniform on the probability simplex.  Their second moments are
$2/[m(m+1)]$ on the diagonal and $1/[m(m+1)]$ off the diagonal, as
follows by dividing independent exponential variables by their sum.
Hence, with $A=UXU^*$ and
$V_X=\operatorname{tr}X^2-(\operatorname{tr}X)^2/m$,
\[
 \operatorname{Var}(A_{11})=\frac{V_X}{m(m+1)},\qquad
 \operatorname{Cov}(A_{11},A_{22})
       =-\frac{V_X}{m(m+1)(m-1)}.
\]
The second identity also follows from permutation symmetry and the
constant value of $\operatorname{tr}A$.  Using the two trace identities
for $D_m$ now gives
\[
 \mathbb E Y^2=\frac{m}{m+1}V_X
                 \le\frac{m^2}{4(m+1)}R(X)^2.
\]
To verify the last inequality, put the entries of $X$ in $[b,b+R]$.
Averaging $(x_j-b)(b+R-x_j)\ge0$ shows that their variance is at most
$(\bar x-b)(b+R-\bar x)\le R^2/4$.

Equation \eqref{eq:unimodal-cosine-bound} proves that
$\mathcal O_m(X)\ge0$ whenever
$R(X)^2\le4\pi^2(m+1)/(3m^2)$.  Taking the infimum proves the first
inequality in \eqref{eq:orbital-unimodal-radius}; the second follows
from $(4/3)(1-m^{-2})\ge1$.  For $m=2$, $D_2$ has entries $-1,1$
and $Y$ is uniform on $[-R(X),R(X)]$, so
$\mathcal O_2(X)=\sin R(X)/R(X)$ and $\mathfrak r_2=\pi$.
\end{proof}

Let $L_{\rm T}=\lim_{a\downarrow0}\sqrt a\,m_*(-a)$ and
$L_{\rm J}=\lim_{a\downarrow0}a\,j_*(-a)$ be the extended limits in
the finite-detection theorems, with $H_0(x)=\Xi(x/2)/8$.  Their formulas
give
\[
 L_{\rm T}=\inf_{\gamma>0:\,m_\gamma\ge2}
             \frac{\gamma\mathfrak r_{m_\gamma}}{\sqrt2},
 \qquad
 L_{\rm J}=\frac{\gamma_*^2}{2},
\]
where $\gamma_*$ is the least positive multiple ordinate; an empty
infimum is $+\infty$.

\begin{corollary}[Finite trace bounds for global detection]
\label{cor:trace-global-threshold}
For every finite compression with $\mathcal D_N>0$,
\[
 L_{\rm T}\ge\frac{\pi}{\sqrt{\mathcal D_N}},
 \qquad L_{\rm J}\ge\frac1{\mathcal D_N}.
\]
When $\mathcal D_N=0$, both limits are infinite.
\end{corollary}

\begin{proof}
The signed trace-defect identity gives
$\mathcal D_N\ge2(m_\gamma-1)/\gamma^2$ at every positive multiple
ordinate.  Therefore
\[
 \frac{\gamma^2\mathfrak r_{m_\gamma}^2}{2}
       \ge\frac{(m_\gamma-1)\mathfrak r_{m_\gamma}^2}{\mathcal D_N}
       \ge\frac{\pi^2}{\mathcal D_N}.
\]
Taking the infimum proves the first assertion.  At the least multiple
ordinate the same defect identity gives
$\gamma_*^2/2\ge(m_{\gamma_*}-1)/\mathcal D_N
\ge1/\mathcal D_N$.  A zero trace defect forces every multiplicity
excess to vanish, and the finite-detection formulas then give both
infinite limits.
\end{proof}
